\begin{abstract}
Deployed content-provenance systems attach a signed manifest to a media file. The manifest is metadata that platforms strip, it identifies the signing device, and its presence is visible. We hide a zero-knowledge proof of origin inside the H.264 bitstream instead. A camera whose key is a leaf of a public Poseidon Merkle registry produces a 129-byte Groth16 proof that it is registered and vouches for a binding of the video and a message; the proof is embedded, under a separate stego key, into CAVLC trailing-one signs of IDR frames, which keeps every codeword length and the file size. To bind the proof to the stream that carries it, we hash every NAL unit with exactly the keyed carrier bits cleared: the digest is invariant under embedding, while any other edit---another video, truncation, re-encoding, an edited non-carrier bit---changes it. Verifiers need the stego key and the registry root but no camera secret, cannot forge proofs and do not learn which camera signed. We also derive a closed-form distortion model: one sign flip costs $4\Qs^2\kappa$ in pixel energy, $\kappa\in[0.925,1.057]$, times a propagation gain $G$, which inverts into a per-frame bit budget for a target PSNR. On 1{,}200 single flips and 810 frames the direct energy matches the decoder exactly, $G$ does not depend on QP, and its heavy-tailed distribution---a tenth of the carriers causes four fifths of the distortion---predicts frame distortion. A 68-byte message and its proof occupy 26 IDRs with a minimum PSNR-Y of 44.03\,dB; proving takes 1.5--1.8\,s, verification about 1\,s, and an 11-case attack matrix including replay is rejected as expected.
\end{abstract}

\begin{IEEEkeywords}
Video steganography, H.264/AVC, CAVLC, zk-SNARK, Groth16, content provenance, anonymous attestation, distortion model.
\end{IEEEkeywords}
